\documentclass{article}
\usepackage{graphicx} % Required for inserting images
\usepackage{xcolor} % For color text like \color{red}
\usepackage{multicol}
\usepackage{xurl}
\renewcommand{\baselinestretch}{1.1}

\setlength {\textwidth}{17.0cm} \setlength {\textheight}{22cm} \setlength {\oddsidemargin}{-0.5cm} \setlength
{\topmargin}{0cm}

\title{Special Session Proposal for IJCNN 2027}
\author{}
\date{}

\begin{document}

\maketitle

\section*{Title}
\noindent
Trustworthy LLMs and AI Agents for Education

\section*{Abstract}
\noindent
Large Language Models (LLMs) and AI agents offer new opportunities to support students and educators through personalized tutoring, learning material generation, automated assessment, and adaptive feedback. Their educational value, however, depends on more than the ability to generate plausible responses. These systems must provide reliable information, accommodate diverse learning needs, protect sensitive data, and support students’ reasoning and autonomy.\\
Educational applications introduce challenges that require dedicated methods and evaluation protocols. Hallucinated explanations and misleading feedback can compromise learning, while biases can affect assessment and the support offered to different student populations. Personalization raises questions about privacy and the appropriate use of learner data. Agentic systems introduce further challenges: errors may propagate across planning, retrieval, memory, and tool use, requiring mechanisms for monitoring, intervention, and teacher oversight. Finally, current evaluation protocols mostly target answer accuracy, failing to assess pedagogical effectiveness, misconception resolution, or independent student reasoning.  \\
This special session aims to bring together researchers developing and evaluating trustworthy LLMs and AI agents for education. We welcome methodological contributions, benchmarks, and empirical studies addressing reliability, fairness, transparency, privacy, and pedagogical effectiveness. \\

Topics of interest include, but are not limited to:
\begin{itemize}
    \item Fabrication (aka Hallucination) detection, mitigation, and uncertainty calibration in educational LLMs.
    \item RAG and knowledge-grounded methods for tutoring and educational content generation.
    \item Evaluation and mitigation of biases and stereotypes in assessment and personalized learning.
    \item Explainability, response verification, and evidence attribution in educational AI.
    %% MAYBE, TO BE CHECKED, IN CASE IT CAN BE DELETED IF WE RUN OUT OF SPACE
    \item Safe planning, memory, tool use, and human oversight in educational AI agents.
    \item Privacy-preserving methods for learner data and personalization.
    \item Benchmarks and cross-model evaluation of reliability, fairness, and pedagogical effectiveness.
    \item Learner autonomy, trust calibration, and mitigation of cognitive over-reliance.
    \item Inclusive and equitable AI for students with SEND and from disadvantaged backgrounds.
    \item Pedagogically relevant assessment of LLMs and AI agents for education.
\end{itemize}

By focusing on the design and evaluation of neural systems for education, the session aligns with IJCNN’s scope in neural networks, artificial intelligence, and intelligent applications. It will provide a focused venue for investigating how generative and agentic AI can deliver educational support that is reliable, inclusive, and pedagogically meaningful.


\section*{Organizers}
\noindent
\begin{itemize}
    \item Allegra Cuzzocrea, University of Naples Federico II , Italy, \texttt{allegra.cuzzocrea@unina.it}
    \item Emilia Lenzi (not confirmed), TU Wien, Austria, \texttt{[MAIL ADDRESS]}
    \item Petia Georgieva, University of Aveiro, Portugal, \texttt{petia@ua.pt}, %https://scholar.google.com/citations?user=aeA48X8AAAAJ&hl=it&oi=ao
    \item Federica Parlapiano, Marche Polytechnic University, Italy, \texttt{f.parlapiano@pm.univpm.it}, \url{https://scholar.google.com/citations?user=Gxo955cAAAAJ&hl=it&oi=ao}, \url{https://federicaparlapiano.github.io/presentation/}
    \item Davide Traini, Marche Polytechnic University, Italy, \texttt{d.traini@pm.univpm.it}, \url{https://scholar.google.com/citations?user=O7ZSzG0AAAAJ&hl=it&oi=ao}, 
    \item Zenglin Xu, Fudan University, China, \texttt{zenglinxu@fudan.edu.cn}
    \item Marcus Liwicki, Luleå University of Technology, Sweden, \texttt{marcus.liwicki@ltu.se}

\end{itemize}

\section*{Primary Contact}
\noindent
Federica Parlapiano, \texttt{f.parlapiano@pm.univpm.it}

\section*{Short Biography}
\noindent

\smallskip
\noindent
\textbf{Petia Georgieva} is Associate Professor with habilitation in Computer Science at the University of Aveiro and a researcher with the Institute of Electronics and Informatics Engineering of Aveiro (IEETA), Portugal. She has been involved in projects devoted to applying AI algorithms for real problems like image processing, 5g/6g network management, material science, brain computer interfaces. Author of peer-reviewed papers, some of which are award-winning, she serves on advisory boards and invited speaker at international conferences. She is a senior member of the International Neural Network Society. She has been a member of the Organising Committee of many international conferences (IbPRIA2027, WCCI-IJCNN 2026, IJCNN2025, IbPRIA 2022, IJCNN 2019, ICANN’07, etc.).

\smallskip
\noindent
\textbf{Federica Parlapiano} is a PhD student in the Department of Information Engineering at Marche Polytechnic University, Italy. Her research focuses on generative AI and complex network analysis, with a particular interest in integrating LLMs and graph-based approaches to build reliable and explainable RAG systems. She actively applies graph theory to the fields of generative AI and Education. She serves as a reviewer for different international journals and conferences, such as  Discover Artificial Intelligence, Information Science, Knowledge-based Systems, and IEEE WCCI 2026. Recently, she served on the PC for the workshop Graphs Across AI: From Structural Reasoning to Foundation Models at IEEE WCCI 2026.


\smallskip
\noindent
\textbf{Davide Traini} is a post-doctoral researcher in the Department of Information Engineering at Marche Polytechnic University, Italy. His research interests include Artificial Intelligence, Graph Neural Networks, Explainable AI, and LLM. He serves as a reviewer for several international journals and conferences, such as ESWA, Discover Artificial Intelligence, AAAI, IEEE IJCNN, and IEEE WCCI. He is a Program Committee member of the 1st GLOW at ACM WWW26, 2nd GLOW at ISWC26, and the 1st GRAAI at IEEE WCCI 2026. He is an organizer of the 2nd Special Session on Graph-Based Solutions for Explainable and Efficient AI at IEEE WCCI 2026.

\smallskip
\noindent
\textbf{Zenglin Xu} is a Distinguished Professor at the Artificial Intelligence Innovation and Incubation (AI 
3
 ) Institute, Fudan University, and an AI Scientist at the Shanghai Academy of AI for Science. He earned his Ph.D. in Computer Science and Engineering from The Chinese University of Hong Kong and completed postdoctoral training at Purdue University and the Max Planck Institute for Informatics. Dr. Xu currently serves as Vice President of the International Neural Network Society (INNS), Editor-in-Chief of Academia AI and Applications, and Senior Action Editor for Neural Networks. He has published over 200 papers in premier venues such as NeurIPS, ICML, ICLR, ACL, and IEEE TPAMI, with more than 19,800 citations (H-index 71). His research centers on trustworthy machine learning, efficient foundation models, and agentic AI, with impactful applications spanning AI for Science, healthcare informatics, finance, and education.

 \smallskip
\noindent
\textbf{Marcus Liwicki} is chaired professor in Machine Learning and vice-rector for AI at Luleå University of Technology. His research interests include machine learning, pattern recognition, artificial intelligence, AI pedagogy, human computer interaction, digital humanities, knowledge management, ubiquitous intuitive input devices, and document analysis. In 2015, at the young age of 32, he received the ICDAR young investigator award, a bi-annual award acknowledging outstanding achievements of in pattern recognition for researchers up to the age of 40. Marcus is Action Editor for Neural Networks, member of many relevant societies (INNS, IAPR, IGS), often in a senior role and in Sweden, especially at LTU, in charge of the strategy for spearheading education with AI.

\section*{Session Webpage}
\noindent
\url{https://}

\section*{Previous Editions}
\noindent


\section*{Committed Papers (tentative)}




\section*{Special Session Details}
\noindent
The proposed Special Session will follow the standard 2-hour format recommended by IJCNN 2027. It will feature five paper presentations (approximately 15 minutes each), followed by a 30-minute panel discussion and open Q\&A with the audience. 

\end{document}
